\documentclass{article}
\usepackage[main,final]{neurips_2026}   % copy neurips_2026.sty from your PA2/PA3 folder
\usepackage[utf8]{inputenc}
\usepackage[T1]{fontenc}
\usepackage{hyperref}
\usepackage{url}
\usepackage{booktabs}
\usepackage{amsfonts}
\usepackage{amsmath}
\usepackage{amssymb}
\usepackage{nicefrac}
\usepackage{microtype}
\usepackage{xcolor}
\usepackage{graphicx}
\usepackage{titlesec}
\usepackage{adjustbox}
\titlespacing*{\section}{0pt}{8pt}{4pt}
\titlespacing*{\subsection}{0pt}{6pt}{3pt}
\titlespacing*{\subsubsection}{0pt}{5pt}{2pt}
\titlespacing*{\paragraph}{0pt}{4pt}{2pt}
\input{tables/notes.tex}   % generated numbers used inside captions
\newcommand{\tab}[1]{\begin{adjustbox}{max width=\linewidth}\input{tables/#1.tex}\end{adjustbox}}

\title{Compressing Weights, Attention Memory and Decoding Steps of Small Qwen2.5 Models}

\author{%
  Zarman Ali Awan \\
  28100092 \\
  \texttt{28100092@lums.edu.pk}
}

\begin{document}

\maketitle
\renewcommand{\arraystretch}{0.8}
\setlength{\textfloatsep}{6pt plus 2pt minus 2pt}
\setlength{\floatsep}{4pt plus 2pt minus 2pt}
\setlength{\intextsep}{4pt plus 2pt minus 2pt}
\setlength{\abovecaptionskip}{3pt}
\setlength{\belowcaptionskip}{2pt}
\setlength{\tabcolsep}{3.5pt}

% =====================================================================================================
\begin{abstract}
We study three ways of making a small language model cheaper on 16\,GB T4 GPUs. \emph{Weight compression:} on
Qwen2.5-1.5B (WikiText-2 perplexity 16.11), Wanda pruning at 50\% sparsity gives 25.5 against 127.8 for magnitude pruning,
and an outlier-based block allocation lowers it by a further 1.06 [0.93, 1.23]; AWQ-style 4-bit quantization reaches 17.7
at 1.07\,GiB (FP16: 2.88\,GiB), and mixed precision at equal packed bytes is no better. Without sparse or low-bit kernels
no compressed model is faster. \emph{Attention and KV memory:} a tiled online-softmax kernel matches FP32 attention
to $4\times10^{-7}$ with 18\,MiB instead of 7.1\,GiB at 8,192 tokens but is about 150$\times$ slower than PyTorch's fused
kernel; at a matched INT4 KV-cache budget, channel-wise-key quantization and question-aware eviction keep 100\% retrieval accuracy while token-wise
keys fall to 0\% and sink+recent eviction to 33\%. \emph{Speculative decoding:} forward-KL LoRA distillation lowers the
draft's KL to the target by 18\% but raises expected acceptance only from 0.862 to 0.867, and every draft is slower than
plain target sampling (0.70--0.81$\times$) because a 0.5B draft forward costs about 84\% of a 1.5B forward at batch size 1.
Code and results: \url{https://github.com/alizarman786-prog/atml-pa4}.
\end{abstract}

% =====================================================================================================
\section{Introduction}
A deployed language model pays for its weights, for attention over a growing context, and for one serial forward pass per
generated token; each cost has well-known reductions, but a nominal saving need not become a practical one. Task~1
compresses Qwen2.5-1.5B~\citep{qwen2024qwen25} by magnitude and Wanda pruning~\citep{sun2024wanda}, OWL-style non-uniform
sparsity~\citep{yin2024owl}, 2:4 sparsity~\citep{mishra2021accelerating} and AWQ-style~\citep{lin2024awq} uniform and
mixed-precision quantization. Task~2 compares exact attention kernels in the spirit of FlashAttention~\citep{dao2022flashattention},
restricted attention on multi-query associative recall (MQAR;~\citealp{arora2024zoology}), and KV-cache
eviction~\citep{xiao2024streamingllm,zhang2023h2o} and quantization~\citep{liu2024kivi} at a matched byte budget. Task~3
distills a draft with off- and on-policy forward KL~\citep{agarwal2024gkd,zhou2024distillspec} for exact speculative
sampling~\citep{leviathan2023fast,chen2023accelerating}. We report changes in absolute units or percentage points (pp)
with bootstrap 95\% CIs, and separate estimated storage from allocated memory and analytical from measured speedups.

\section{Experimental Setup}
All experiments ran on Kaggle with two Tesla T4 GPUs (16\,GB each; one experiment stream per GPU), PyTorch~2.11.0
(CUDA~13.0)~\citep{paszke2019pytorch}, Transformers~4.57.1~\citep{wolf2020transformers} and PEFT~0.17.1~\citep{peft}, with seed 6304
everywhere. Models are Qwen2.5-1.5B (revision \texttt{8faed76}), Qwen2.5-1.5B-Instruct (\texttt{989aa79}) and
Qwen2.5-0.5B-Instruct (\texttt{7ae5576}) in FP16. Data come locked from the course repository (C4~\citep{raffel2020c4},
WikiText-2~\citep{merity2017wikitext}, HellaSwag~\citep{zellers2019hellaswag}, GSM8K~\citep{cobbe2021gsm8k}, MQAR, retrieval
prompts, prefix banks); every table and figure is generated from saved results by \nolinkurl{report/make_report_assets.py}.

% =====================================================================================================
\section{Task 1: Pruning and Quantization of Qwen2.5-1.5B}
\label{sec:task1}
\subsection{Setup and Design Choices}
All seven linear projections of the 28 blocks are compressed. Calibration uses 128 C4 sequences of 512 tokens, model selection uses a separate development set of 32 C4 sequences
(NLL), and evaluation uses WikiText-2 and C4-validation perplexity over \TOneWikiN{} and \TOneCFourN{} scored tokens in
fixed windows, plus zero-shot HellaSwag acc$_{\text{norm}}$ (lm-eval-harness conventions) on all 2,000 examples for the
central conditions and on the first 512 for the rest ($\dagger$). P1 prunes by $|W_{ij}|$ and P2 by the Wanda score
$|W_{ij}|\,\|X_j\|_2$, both per output row, at $s\in\{0.25,0.5,0.75\}$; P2 is also run with a 2:4 pattern. P3 allocates
block sparsities $s_\ell=S+\lambda(1-2r_\ell)$ from a min-max-normalized sensitivity $r_\ell$, clipped to [0.10, 0.85] and
shifted so that the exact number of zeros matches the uniform model. Two sensitivity measures are candidates: the OWL
outlier ratio (share of Wanda scores above $5\times$ the block mean) and the NLL increase when a single block is pruned
at $S$ (16 calibration sequences); the measure and $\lambda\in\{0.04,0.08,0.12,0.16\}$ are chosen on development NLL.
Q1 is round-to-nearest and Q2 is AWQ-style quantization, both with groups of 128, asymmetric codes, an FP16 scale and a
$b$-bit zero point; Q2 searches an exponent $\alpha\in\{0,0.1,\dots,1\}$ per block and scaling group (qkv, o, gate+up,
down) and folds the scales into RMSNorm, \texttt{v\_proj} and \texttt{up\_proj}, which leaves the FP16 function unchanged.
Q3 assigns 2, 3, 4 or 8 bits to each of the \TOneQThreeUnits{} (block, scaling group) units by an exact multiple-choice
knapsack under the packed-byte budget of uniform Q2 at 3 or 4 bits. Unit costs come from one of three probes (isolated
NLL increase, mean relative output error, or NLL change around the uniform operating point), chosen on development NLL.
Storage counts packed codes plus metadata (checked by real bit-packing) or FP16 nonzeros plus a 1-bit mask;
evaluation runs all weights as dense FP16 tensors.

\subsection{Results}
Table~\ref{tab:t1main} lists all conditions, Table~\ref{tab:t1paired} the paired adaptive-minus-uniform differences and
Fig.~\ref{fig:t1} quality against sparsity and size; selection, allocations and costs are in Tables~\ref{tab:t1sel},
\ref{tab:t1cost} and Fig.~\ref{fig:t1alloc}.

\begin{table}[!tb]
\centering\footnotesize
\caption{Task 1. All conditions: achieved sparsity or effective bits per weight (with metadata), total model size,
WikiText-2 and C4 perplexity with the change from FP16 and a 95\% bootstrap CI over evaluation windows, HellaSwag
acc$_{\text{norm}}$ (\%; $\dagger$: first 512 examples, on which even the 8-bit models score only 58\%, so these values
are not comparable with the 2,000-example ones), and development NLL.}
\label{tab:t1main}
\tab{t1_main}
\end{table}

\begin{table}[!tb]
\centering\footnotesize
\caption{Task 1. Paired differences, adaptive minus uniform rule at the same sparsity or packed-byte budget, with 95\%
bootstrap CIs (perplexity: whole evaluation windows resampled and pooled NLL recomputed; HellaSwag: examples resampled).
Negative perplexity and positive accuracy differences favour the adaptive rule.}
\label{tab:t1paired}
\tab{t1_paired}
\end{table}

\begin{figure}[!tb]
\centering
\includegraphics[width=\linewidth]{figures/fig_task1.pdf}
\caption{Task 1. (a) WikiText-2 perplexity vs.\ target sparsity (log scale; dashed: FP16). (b) HellaSwag
acc$_{\text{norm}}$ (hollow markers: 512-example subset; dotted: chance, 25\%). (c) Perplexity vs.\ packed model size for
Q1, Q2 and Q3.}
\label{fig:t1}
\end{figure}

\subsection{Analysis}
\paragraph{Magnitude, activation and removal impact (RQ1).}
Magnitude and Wanda keep mostly the same weights: at 50\% sparsity the kept sets overlap by 0.86--0.94 per projection,
least for \texttt{o\_proj} and \texttt{q\_proj} (Table~\ref{tab:t1scores}). The disagreement sits where input activations
are uneven: the largest input-feature norm is typically 9--24$\times$ the median and up to 2,816$\times$ for
\texttt{down\_proj}, so a small weight on an outlier feature can carry a large contribution that magnitude pruning
removes. This small set decides quality. At 25\% sparsity magnitude pruning costs +5.22 [4.75, 5.80] WikiText-2
perplexity against +0.44 [0.37, 0.52] for Wanda; at 50\%, +111.7 against +9.39, with HellaSwag at 40.1\% against 58.1\%
(FP16: 69.3\%). At 75\% both collapse ($1.6\times10^6$ and $6.4\times10^4$), so at high sparsity even an activation-aware
per-row score removes weights whose joint effect matters. The 2:4 pattern costs +51.9 instead of +9.4 at the same 50\%,
because it forces zeros into groups where Wanda would have kept several weights.

\paragraph{Sensitivity and allocation (RQ2).}
P3 selected the OWL outlier ratio at both budgets (Table~\ref{tab:t1sel}). At $S=0.5$ the best $\lambda$ is small (0.08;
block sparsities 0.38--0.54) and larger values hurt; at $S=0.75$ development NLL falls from 10.30 to 5.38 and still
improves at the largest $\lambda$ tested (blocks 0.52--0.84). A uniform 75\% budget therefore destroys a few outlier-heavy
blocks, while at 50\% only mild reallocation helps. For Q3 the in-context probe wins at both budgets and relative output
error is worst (+0.11 and +0.19 development NLL), since a projection's own output error does not say how much the loss
depends on it. Even the best allocation does not beat uniform AWQ on the development set (+0.013 at 4 bits, +0.002 at
3 bits): it raises 11 units to 8 bits by lowering 15 to 3 bits, and its additive per-unit costs miss the interactions
that make several 3-bit units together worse than predicted.

\paragraph{Uniform vs.\ layer-dependent budgets (RQ3).}
At matched budgets (sparsity exactly 0.500 and 0.750; Q3 at 4.15 and 3.15 effective bits vs.\ 4.16 and 3.15), the paired
differences support the development selection only for pruning (Table~\ref{tab:t1paired}). P3 lowers perplexity by
1.06 [0.93, 1.23] (WikiText-2) and 0.50 (C4) at 50\%, yet its HellaSwag accuracy is $-0.8$ [$-1.6$, +0.1]\,pp; at 75\%
perplexity falls from 63,771 to 282 while HellaSwag moves only +3.5 [$-1.0$, +8.0]\,pp, because both models are near
chance and perplexity, an exponential of mean NLL, magnifies differences between broken models. Q3 is worse than Q2 at
4 bits (+0.24 and +0.22, CIs excluding zero) and indistinguishable at 3 bits, with HellaSwag differences inside their CIs.
HellaSwag depends on the ranking of four endings, which a change in which blocks are damaged can flip without moving
mean token loss (example 18 in Table~\ref{tab:t1qual}: P2 and Q2 keep FP16's answer, P3 and Q3 do not); with CI
half-widths of about 1\,pp, perplexity is the more sensitive measure.

\paragraph{Theoretical vs.\ practical savings (RQ4).}
Calibration error tracks quality between methods (P1 vs.\ P2, Q1 vs.\ Q2) but not between allocations: P3 has a higher
mean output error than P2 at 50\% (0.0286 vs.\ 0.0275) yet lower perplexity (Table~\ref{tab:t1cost}). Applying a
compression takes 0.3--6\,s, but the calibration pass (\TOneCalibS\,s), AWQ search (\TOneAwqS\,s per bit width) and P3/Q3
probes (\TOnePThreeProbeMin{} and \TOneQThreeProbeMin{} min) are one-time costs that only the adaptive rules need in full.
Storage falls less than nominally: 50\% sparsity saves 37\% (1.81\,GiB) because the bit-mask and FP16 embeddings remain,
4-bit packing with metadata (4.16 bits/weight) saves 63\%. Latency does not fall: prefill of 512 and 2,048 tokens stays at
141--155 and 892--943\,ms and 64 decode steps at about 2\,s, because zeros and dequantized codes run as dense FP16
products; the T4 has no sparse tensor cores for 2:4 and we used no low-bit kernels. Outliers up to $1.7\times$ (P2 2:4)
have identical arithmetic and reflect contention from the concurrent Task~2 stream.

% =====================================================================================================
\section{Task 2: Efficient Attention and Context Memory}
\label{sec:task2}
\subsection{Setup and Design Choices}
\emph{Part A} compares dense attention, an online-softmax tiled kernel (128$\times$128 tiles, FP32 running statistics,
backward by autograd through the Python loop) and PyTorch's fused scaled-dot-product attention (SDPA) for causal
attention with \TTwoAShape, $L=256$--8,192, FP16 and FP32, against an FP32 dense reference; timings are medians of
\TTwoABenchReps, memory is peak above the inputs. \emph{Part B} trains 2-layer MQAR models (width 128, four 32-dim heads,
$T=256$, 4,500 AdamW updates of batch 32, \TTwoBSeeds{} seeds) with dense, window-64 (126-token reach over two layers),
local+dilated (two window-64 and two stride-4 heads; reach 504), linear (ELU+1; 4,224 state floats per layer) or hybrid
attention, on near (distance $<64$) or mixed distances, tested on near, far ($\ge128$) and mixed splits for
$K\in\{4,8,16\}$. The hybrid intervention (two linear and two local-16 softmax heads; 4,160 state floats, matched to linear)
was motivated by local heads retrieving near pairs exactly while linear heads summarize distant ones. Chance is $1/K$
among present values and 1/64 over labels. \emph{Part C} runs Qwen2.5-1.5B-Instruct on
the \TTwoCN{} staff retrieval prompts (2,048 and 4,096 tokens; evidence at start, middle or end; 30 per cell) with greedy
decoding of up to 16 tokens and exact match of the code. The budget is the whole cache in INT4 (groups of 64, FP16 scale
and offset), 0.28--0.29$\times$ FP16. \emph{Sink+recent} keeps 8 sinks and the most recent tokens; \emph{question-aware}
eviction keeps 8 sinks, the 64 most recent tokens and the positions receiving most attention from the question
(summed over heads, max-pooled), counting index bytes; \emph{token-wise INT4} groups keys and values over 64 features of a
token, \emph{KIVI-style INT4} groups keys over 64 tokens of a channel and values per token; the \emph{hybrid} keeps
question-selected tokens in FP16 and the rest in INT2 (KIVI layout). Evicted caches keep original RoPE positions; quality
uses dequantized caches, bytes are checked by real packing.

\subsection{Results}
Tables~\ref{tab:t2abench}, \ref{tab:t2b} and \ref{tab:t2c} with Figs.~\ref{fig:t2b} and~\ref{fig:t2c} give the main results;
latency curves, fidelity, tile sizes, theory, controls, the inference sweep and packing costs are in the appendix.

\begin{table}[!tb]
\centering\footnotesize
\caption{Task 2A. FP16 median latency (ms) and peak memory above the inputs (MiB) for the forward and the forward+backward
pass at three sequence lengths (\TTwoAShape); OOM: out of memory on a 16\,GB T4.}
\label{tab:t2abench}
\tab{t2a_bench}
\end{table}

\begin{table}[!tb]
\centering\footnotesize
\caption{Task 2B. MQAR test accuracy (\%, mean $\pm$ s.d.\ over \TTwoBSeeds{} seeds) for
each trained configuration on six of the nine test splits; chance is $100/K$ for value-aware guessing and 1.6 over the
label space. Train: minutes per run.}
\label{tab:t2b}
\tab{t2b_mqar}
\end{table}

\begin{figure}[!tb]
\centering
\includegraphics[width=\linewidth]{figures/fig_task2b.pdf}
\caption{Task 2B. (a) K16\_mixed test accuracy by key--query distance for the models trained at $K=16$ on mixed
distances; (b) their development accuracy during training (dotted: $1/K$); (c) single-layer forward time against $T$
(batch 8, FP16).}
\label{fig:t2b}
\end{figure}

\begin{table}[!tb]
\centering\footnotesize
\caption{Task 2C. Retrieval accuracy (\%, 95\% prompt-bootstrap CI, 90 prompts per length), stored KV bytes relative to
FP16, share of prompts within budget, share of evidence tokens retained, first-token KL and mean reference-answer NLL
relative to the full cache, selection+compression time and peak allocated memory.}
\label{tab:t2c}
\tab{t2c_policies}
\end{table}

\begin{figure}[!tb]
\centering
\includegraphics[width=\linewidth]{figures/fig_task2c.pdf}
\caption{Task 2C. (a) Accuracy against context length (95\% CI); (b) reference-answer NLL against stored bytes relative
to FP16; (c) accuracy by evidence position.}
\label{fig:t2c}
\end{figure}

\subsection{Analysis}
\paragraph{Exactness is cheap; speed needs a fused kernel (RQ1).}
All implementations are numerically sound (Table~\ref{tab:t2afid}): in FP32 the tiled kernel is within $3.6\times10^{-7}$
of the reference in the output and $1.1\times10^{-6}$ in the gradients, and in FP16 its output error
($3.9$--$4.0\times10^{-4}$) is the smallest, below SDPA ($4.2\times10^{-4}$) and dense ($5.8\times10^{-4}$), because it keeps
the softmax statistics in FP32; tile size changes neither error nor memory much (8--11\,MiB; Table~\ref{tab:t2atiles}).
Memory follows the algorithm (Table~\ref{tab:t2abench}, Fig.~\ref{fig:t2a}): at 8,192 tokens the dense forward allocates
7,240\,MiB, the tiled forward 18\,MiB and SDPA 8\,MiB; dense forward+backward runs out of memory, the tiled backward needs
3,713\,MiB because autograd stores every tile's probabilities, and SDPA's fused backward recomputes them in 33\,MiB. Speed
follows the implementation: launching several small kernels per tile from Python makes the tiled kernel 4--10$\times$ slower
than dense and about 150$\times$ slower than SDPA in the forward pass (709 vs.\ 4.7\,ms), 121$\times$ with backward. Online
softmax removes the $L^2$ buffer; the speed of FlashAttention comes from fusing the loop.

\paragraph{Reach, selection and state in associative recall (RQ2).}
Architecture limits what can be represented, but at $K=16$ training decides what is learned (Table~\ref{tab:t2b}). Dense
attention trained on near pairs reaches 99.8--100\% on every split, far ones included, while the same model trained on
mixed distances ends at 6.5\% ($\approx1/K$) after 4,500 updates (Fig.~\ref{fig:t2b}b), like every other model trained at
$K=16$ on mixed distances; dense models trained at $K=4$ and 8 solve their own $K$ and reach 32.7\% and 53.0\% at $K=16$.
The window model shows the access limit: trained on near pairs it scores 100\% on near, 22.5\% on mixed and 1.6\%
(label-space chance) on far queries, which lie beyond its 126-token reach. A dense near-trained model given a window-64
mask at test time reproduces exactly these numbers, whereas keeping only its top-64 keys per query leaves it at 100\%
(Table~\ref{tab:t2bctl}): recall needs long-range access to only a few positions, so content selection suffices if it picks
the right keys. Dilation reaches every distance (504) but did not learn ($\le$28.5\% even at $K=4$); linear attention at
$K=4$ varies strongly between seeds (44.5 $\pm$ 31.4\%); and the hybrid did no better than linear attention (6.2\% at
K16\_mixed and on near queries), contrary to its motivation. Since dense attention fails in the same mixed $K=16$ setting,
we read these failures as optimization failures under this budget rather than capacity limits.

\paragraph{Eviction, quantization and hybrids at the INT4 budget (RQ3).}
At 0.28--0.29$\times$ the FP16 bytes, KIVI-style quantization and question-aware eviction keep 100\% accuracy, the hybrid
98.9--100\%, sink+recent 33.3\% and token-wise quantization 0\% (Table~\ref{tab:t2c}). Position explains sink+recent
completely: it answers every prompt with evidence at the end and none of the others (Table~\ref{tab:t2cpos}). The two INT4
layouts differ only in key grouping, which alone moves accuracy from 0\% to 100\%: keys contain a few channels of much
larger magnitude, which set the step size of a group spanning 64 channels of one token and collapse the other channels
to a few levels, whereas per-channel groups give each channel its own range, the motivation of KIVI~\citep{liu2024kivi}.
The token-wise cache shifts the first answer token by a KL of 5.5--5.6 nats and yields no code (Table~\ref{tab:t2cqual}).
Among the policies at 100\%, likelihood still separates them: question-aware eviction (first-token KL 0.001--0.002) and
KIVI (0.002--0.003) are nearly lossless, while the hybrid's INT2 tokens perturb more (0.006--0.016) and cost one answer;
Fig.~\ref{fig:t2c}b shows the same ordering in answer NLL. Evidence retention explains eviction: question-aware selection
keeps 96--97\% of evidence tokens, sink+recent 33\%. Question-aware eviction is optimistic here, as it sees the question
before compressing.

\paragraph{Overheads behind nominal compression (RQ4).}
Real packing reproduces the byte formula exactly (15.8 vs.\ 56\,MiB at batch 1 and 2,048 tokens; Table~\ref{tab:t2cpack}),
but the saving becomes neither speed nor peak memory: packing a batch-8, 4,096-token cache takes 177--183\,ms and unpacking
149\,ms, a decode step with the packed cache is 31--73\% slower (45.7 vs.\ 34.9 and 381 vs.\ 221\,ms) because it is
dequantized to FP16 before SDPA every step, and peak allocated memory rises slightly (5.20 vs.\ 5.07\,GiB at batch 8)
because weights and activations dominate. Question scoring makes question-aware eviction take 74--106\,ms and the hybrid
98--142\,ms, against 17--34\,ms for quantization and 8--17\,ms for sink+recent. In Part~B the variants with fewer attended
pairs (14,368 for window-64 vs.\ 32,896 for dense at $T=256$; Table~\ref{tab:t2btheory}) run 2--5$\times$ slower than dense
SDPA (Table~\ref{tab:t2bsweep}), because they pass a materialized boolean mask (dilated: 1,264\,MiB at 8,192) or a chunked
FP32 scan.

% =====================================================================================================
\section{Task 3: Policy Distillation and Speculative Decoding}
\label{sec:task3}
\subsection{Setup and Design Choices}
The target is Qwen2.5-1.5B-Instruct and the draft Qwen2.5-0.5B-Instruct (shared tokenizer), sampling at temperature 1
throughout. Our exact speculative sampler accepts with probability $\min(1,p/q)$, resamples from the normalized residual
$(p-q)_+$, adds a bonus token after $\gamma$ acceptances and rolls back both KV caches; it passes \TThreeUnitTests{} unit
tests and a $\chi^2$ test on a toy Markov model with a wrong-residual negative control (Table~\ref{tab:t3toy}). Distillation
trains LoRA~\citep{hu2022lora} adapters ($r=8$, $\alpha=16$ on q/k/v/o; \TThreeTrainable{} parameters) with token-level
forward KL(target$\,\|\,$draft) on responses to 500 GSM8K-train and 300 short-form training prompts, sampled from the target
(off-policy) or the current student (on-policy) with caps of 192 and 32 tokens; each run uses 30k math and 10k
short-form scored tokens in exact 256-token updates (157 updates, learning rate $3\times10^{-4}$), two seeds, and merged
adapters. Alignment (FKL, RKL, TV, argmax agreement, $1-$TV) uses \TThreeAlignPrefixes{} prefixes per draft from seven
locked banks. The speculative benchmark runs the 120 test prompts (60 GSM8K, 30 familiar and 30 shifted short-form) for up
to 64 tokens with $\gamma\in\{2,4\}$ and \TThreeSpecReps{} counterbalanced repetitions, each paired with plain target
sampling under the same seed; speedup is the geometric mean of tokens-per-second ratios.

\subsection{Results}
Table~\ref{tab:t3align} gives alignment, Table~\ref{tab:t3spec} and Fig.~\ref{fig:t3} speculative decoding; training costs,
standalone accuracy and diagnostics are in the appendix (Tables~\ref{tab:t3toy}--\ref{tab:t3acc}).

\begin{table}[!tb]
\centering\footnotesize
\caption{Task 3. Expected one-token acceptance $1-$TV between target and draft, averaged over the prefixes of each bank:
target samples, original-draft samples, exact speculative output ($\gamma=4$), and each student's own rollouts.
FKL in Table~\ref{tab:t3fkl}; all metrics with CIs in Table~\ref{tab:t3alignfull}.}
\label{tab:t3align}
\tab{t3_align_accept}
\end{table}

\begin{table}[!tb]
\centering\footnotesize
\caption{Task 3. Speculative decoding per draft and $\gamma$ (\TThreeSpecPrompts{} prompts $\times$ \TThreeSpecReps{}
repetitions, up to 64 tokens): acceptance rate, accepted and emitted tokens per round, tokens per round predicted by
$(1-\alpha^{\gamma+1})/(1-\alpha)$ with $\alpha$ the mean $1-$TV on target prefixes, tokens per target call, throughput,
and geometric-mean speedup over plain target sampling with a 95\% bootstrap CI.}
\label{tab:t3spec}
\tab{t3_spec}
\end{table}

\begin{figure}[!tb]
\centering
\includegraphics[width=\linewidth]{figures/fig_task3.pdf}
\caption{Task 3. (a) Forward KL per scored token during distillation (rolling mean over 10 updates); (b) predicted
against measured tokens per round (dotted: identity); (c) measured speedup against $\gamma$ (95\% CI; dotted: no speedup).}
\label{fig:t3}
\end{figure}

\subsection{Analysis}
\paragraph{Prefix distribution and agreement (RQ1).}
Distillation reduced the quantity it optimizes far more than the one speculative decoding uses: on target prefixes
forward KL falls from 0.216 to 0.177 for both off-policy seeds (0.179--0.183 on-policy), but $1-$TV moves only from 0.862
to 0.867 with overlapping CIs (Tables~\ref{tab:t3align}, \ref{tab:t3alignfull}). Off-policy drafts are at least as well
aligned as on-policy ones on every bank, including the students' own rollouts (off-policy seed~1: 0.871 vs.\ 0.859 for the
original draft on the same prefixes; on-policy seed~1: 0.845 vs.\ 0.841). On-policy training also raised reverse KL on
target prefixes (0.261 and 0.276 vs.\ 0.240) while lowering forward KL, consistent with mass-covering forward-KL
training. Prefixes written by drafts, the original's and the on-policy students' own, are the hardest for every draft
(0.826--0.850), so on-policy training spent its 157 updates on unusual text. The sampler is exact in the toy test (correct: $p=0.27$--$0.95$;
wrong residual: $\chi^2>12{,}900$), yet the exact-speculation bank shows higher agreement than the target bank (0.881
[0.875, 0.887] vs.\ 0.862 [0.857, 0.868]); as both follow the target distribution, we attribute this to which prefixes were
drawn: each bank holds one rollout per prompt, and the prefix bootstrap ignores correlation within a rollout.

\paragraph{Domain, length and transfer (RQ2).}
Distillation helped only where the training data came from (Table~\ref{tab:t3aligndom}): on familiar short-form prompts
agreement rises from 0.772 to 0.818--0.827 and acceptance at $\gamma=4$ from 55.0\% to 60.6--63.0\% (Table~\ref{tab:t3specdom}),
but on the shifted split it stays at 0.799--0.811 (original 0.814), and on mathematics, three quarters of the training
tokens, at 0.859--0.868, where the original draft already agreed well. Response length dominates speed: math runs reach
the 64-token cap and achieve 0.90--0.92 at $\gamma=4$, short-form answers stop after 12--23 tokens and reach 0.51--0.58
because two prefills and per-call overhead are spread over few tokens. Acceptance per proposal falls from about 82\% at
$\gamma=2$ to 68\% at $\gamma=4$, but mostly mechanically: with i.i.d.\ acceptance $\alpha=0.862$ a round accepts
$\alpha(1-\alpha^\gamma)/(1-\alpha)$ tokens, i.e.\ 80\% and 70\% of the proposals.

\paragraph{From KL and TV to tokens and time (RQ3).}
$1-$TV on target prefixes predicts the work per round well: predicted tokens per round (2.59--2.62 at $\gamma=2$,
3.76--3.84 at $\gamma=4$) exceed the measured ones by only 2--5\% (Fig.~\ref{fig:t3}b), mainly because the cap and EOS cut
the last round short. KL is a poor predictor: 18\% lower forward KL bought 0.005 in $1-$TV and no measurable speed. Measured
costs explain the speed (Table~\ref{tab:t3decode}): a draft forward takes 27--30\,ms against 34--37\,ms for the target,
and verifying 3 or 5 tokens costs no more than one. At batch size 1 the T4 is bound by per-layer overhead, so cost follows
the number of layers (24 vs.\ 28) rather than parameters (0.5B vs.\ 1.5B). With $c_D\approx28.5$\,ms, $c_T\approx34$\,ms
and $\tau$ tokens per target call, the analytical speedup $c_T\tau/(\gamma c_D+c_T)$ is 0.90 ($\gamma=2$) and 0.79
($\gamma=4$) for the original draft against 0.80 and 0.71 measured; in the fixed-64-token math diagnostic
(Table~\ref{tab:t3fixed}) it is 0.95 and 0.90 against 0.96 and 0.92 measured, so per-round host overhead is small and the
remaining gap comes from short-form prompts where prefill dominates.

\paragraph{When the costs dominate (RQ4).}
The draft's own cost dominates: with $\tau\approx2.6$ ($\gamma=2$) and $\approx3.8$ ($\gamma=4$), speculative decoding
breaks even only if a draft forward costs less than about 0.8 or 0.7 of a target forward, while ours costs about 0.84, so
no draft and no $\gamma$ beats the target (0.70--0.81$\times$, all CIs below 1; Fig.~\ref{fig:t3}c). Unmerged adapters
would be worse still (18.8 vs.\ 26.4 tokens/s; Table~\ref{tab:t3merge}), and EOS termination hurts mainly short prompts.
Distillation adds cost (Table~\ref{tab:t3cost}): off-policy 2.8 minutes of teacher generation plus 2.7 of
training, on-policy 32--34 minutes, 30 of them rollouts, six times more for no better alignment, and neither is recovered
without a speedup. Distillation also changed
the drafts' own behaviour (Table~\ref{tab:t3acc}): GSM8K accuracy falls from 25.0\% to 0.8--10.8\% and truncation at 512
tokens rises from 39\% to 77--88\%, likely because math training responses were capped at 192 tokens, so the students
rarely saw an end-of-sequence token after long reasoning (Table~\ref{tab:t3qual}). Speculative outputs are unaffected,
since they follow the target distribution exactly.

% =====================================================================================================
\section{Conclusion}
\label{sec:conclusion}
In all three tasks the nominal saving was easy and the practical one was not. Activation-aware pruning and AWQ-style
scaling preserved quality far better than naive rules, but without sparse or low-bit kernels no compressed model was
faster. Tiling removed the quadratic buffer but was slow without fusion; at the INT4 cache budget key grouping mattered
more than eviction versus quantization, and packing added decode time. Distillation lowered the draft's KL, but the draft
cost nearly as much per forward as the target, so speculative decoding lost. The deciding quantities were paired
differences, evidence retention, $1-$TV and measured forward costs rather than nominal ratios.

\paragraph{Limitations and deviations.} Tasks~1 and 2C use one seed (MQAR three, distillation two) and some HellaSwag values
use the 512-example subset; Task~1 latencies were measured while Task~2 ran on the other GPU. The course prefix-bank script
raised an error with the pinned Transformers~4.57.1 (\texttt{generate} rejects \texttt{generator=}) and was run unchanged
through a wrapper seeding the RNG from the same per-prompt seed. The first Task~3 alignment run exhausted host memory and
was recomputed separately from the saved checkpoints. Dense forward+backward attention ran out of memory at 8,192 tokens and
the tiled kernel is an unfused reference. The MQAR hybrid's motivation was
not recorded before evaluation, and question-aware eviction assumes the question is known before compression.

\bibliographystyle{plainnat}
\bibliography{references}

% =====================================================================================================
\clearpage
\appendix
\setlength{\floatsep}{14pt plus 2pt minus 2pt}\setlength{\intextsep}{14pt plus 2pt minus 2pt}
\section{Additional results}

\begin{table}[!htb]
\centering\footnotesize
\caption{Task 1. Compression cost and latency (median ms, batch 1, dense FP16 execution; prefill of 512 and 2,048 tokens,
64 decode steps after a 512-token prefix), storage format and mean relative output error on the calibration tokens.
Compress time excludes the one-time calibration pass (\TOneCalibS\,s), the AWQ search (\TOneAwqS\,s per bit width) and
the P3/Q3 probes (\TOnePThreeProbeMin{} and \TOneQThreeProbeMin{} min).}
\label{tab:t1cost}
\tab{t1_cost}
\end{table}

\begin{table}[!htb]
\centering\footnotesize
\caption{Task 1. P3 and Q3 selection on the 32-sequence development set: every candidate rule, its sensitivity measure,
$\lambda$, the resulting block-sparsity range (P3) or number of units at 2/3/4/8 bits (Q3), development NLL and the
change from the uniform rule at the same budget; \checkmark: selected. Q3 selects among non-uniform candidates only.}
\label{tab:t1sel}
\tab{t1_selection}
\end{table}

\begin{figure}[!htb]
\centering
\includegraphics[width=\linewidth]{figures/fig_task1_alloc.pdf}
\caption{Task 1. (a) Selected P3 sparsity per block for $S=0.5$ and $0.75$ (dotted: uniform). (b, c) Selected Q3 bit
widths per block and scaling group at the 3- and 4-bit budgets.}
\label{fig:t1alloc}
\end{figure}

\begin{table}[!htb]
\centering\footnotesize
\caption{Task 1. Per-projection score statistics over the 28 blocks: kept-set overlap between magnitude and Wanda pruning
at 50\% and 75\% sparsity, ratio of the largest to the median input-feature norm, and share of Wanda scores above
5$\times$ the mean.}
\label{tab:t1scores}
\tab{t1_scores}
\end{table}

\begin{figure}[!htb]
\centering
\includegraphics[width=\linewidth]{figures/fig_task1_scores.pdf}
\caption{Task 1. (a) Magnitude/Wanda kept-set overlap at 50\% sparsity per block and projection; (b) activation-norm
outlier ratio per block.}
\label{fig:t1scores}
\end{figure}

\begin{table}[!htb]
\centering\footnotesize
\caption{Task 2A. Numerical fidelity against the FP32 dense reference (maximum over sequence lengths): relative
Frobenius error of the output and of the three gradients, output max-abs error, and whether all outputs are finite.}
\label{tab:t2afid}
\tab{t2a_fidelity}
\end{table}

\begin{figure}[!htb]
\centering
\includegraphics[width=\linewidth]{figures/fig_task2a.pdf}
\caption{Task 2A. Median latency of (a) the forward and (b) the forward+backward pass, and (c) peak memory above the
inputs for forward+backward, against sequence length (FP16, log-log).}
\label{fig:t2a}
\end{figure}

\begin{table}[!htb]
\centering\footnotesize
\caption{Task 2A. Tile-size sweep of the tiled kernel at $L=2{,}048$ (FP16 forward): error against FP32, extra peak memory,
and size of one FP32 score tile (the full score matrix would need 128\,MiB).}
\label{tab:t2atiles}
\tab{t2a_tiles}
\end{table}

\begin{table}[!htb]
\centering\footnotesize
\caption{Task 2B. Theoretical attended pairs per head at $T=256$, reachable distance after two layers, and decode-time
state per layer in floats ($t$: current length).}
\label{tab:t2btheory}
\tab{t2b_theory}
\end{table}

\begin{table}[!htb]
\centering\footnotesize
\caption{Task 2B. Dense models evaluated as trained, with a test-time window-64 mask, or with only the top-64 keys per
query (selected by the dense scores); accuracy in \%.}
\label{tab:t2bctl}
\tab{t2b_controls}
\end{table}

\begin{table}[!htb]
\centering\footnotesize
\caption{Task 2B. Forward-only inference sweep of one attention layer with random inputs (batch 8, FP16): median time in
ms / extra peak memory in MiB. Window and dilated attention use SDPA with a materialized boolean mask; linear attention
uses a chunked (64) FP32 scan.}
\label{tab:t2bsweep}
\tab{t2b_sweep}
\end{table}

\begin{table}[!htb]
\centering\footnotesize
\caption{Task 2C. Accuracy (\%) by evidence position, pooled over both context lengths (60 prompts per position).}
\label{tab:t2cpos}
\tab{t2c_position}
\end{table}

\begin{table}[!htb]
\centering\footnotesize
\caption{Task 2C. Real packed INT4 cache (KIVI and token layouts): FP16, serialized packed and formula bytes,
pack/unpack time, decode-step time with the FP16 and the packed cache, and peak allocated memory. The packed path
dequantizes to FP16 before SDPA.}
\label{tab:t2cpack}
\tab{t2c_packing}
\end{table}

\begin{table}[!htb]
\centering\footnotesize
\caption{Task 3. Exactness test of the speculative sampler: $\chi^2$ against the exact target sequence distribution of a
toy Markov model for the correct sampler and for a negative control that resamples from $p$ instead of the residual.}
\label{tab:t3toy}
\tab{t3_toy}
\end{table}

\begin{table}[!htb]
\centering\footnotesize
\caption{Task 3. Forward KL(target $\|$ draft) per bank, mean over prefixes.}
\label{tab:t3fkl}
\tab{t3_align_fkl}
\end{table}

\begin{table}[!htb]
\centering\footnotesize
\caption{Task 3. All alignment metrics per draft and bank with 95\% prefix-bootstrap CIs for FKL and $1-$TV
(\TThreeAlignRows{} draft--prefix rows).}
\label{tab:t3alignfull}
\tab{t3_align_full}
\end{table}

\begin{table}[!htb]
\centering\footnotesize
\caption{Task 3. Mean $1-$TV by prompt domain over all banks.}
\label{tab:t3aligndom}
\tab{t3_align_domain}
\end{table}

\begin{table}[!htb]
\centering\footnotesize
\caption{Task 3. Speculative decoding by prompt domain at $\gamma=\TThreeSpecGammaMax$: acceptance, mean emitted tokens,
tokens per target call and geometric-mean speedup.}
\label{tab:t3specdom}
\tab{t3_spec_domain}
\end{table}

\begin{table}[!htb]
\centering\footnotesize
\caption{Task 3. Fixed-length diagnostic: 24 GSM8K prompts, exactly 64 tokens (EOS ignored), one repetition; mean wall
time of speculative and plain target decoding and geometric-mean speedup.}
\label{tab:t3fixed}
\tab{t3_fixed64}
\end{table}

\begin{table}[!htb]
\centering\footnotesize
\caption{Task 3. Per-forward decode cost (median ms, batch 1) of the target and each merged draft for 1, 3 and 5 new
tokens after a prefix of 64, 256 or 512 tokens.}
\label{tab:t3decode}
\tab{t3_decode}
\end{table}

\begin{table}[!htb]
\centering\footnotesize
\caption{Task 3. Generation of 64 tokens with the LoRA adapter merged into the weights or applied unmerged.}
\label{tab:t3merge}
\tab{t3_merge}
\end{table}

\begin{table}[!htb]
\centering\footnotesize
\caption{Task 3. Distillation runs: optimizer updates, updates skipped (non-finite loss), scored tokens, training
sequences, cost breakdown in minutes (teacher generation, student rollouts, teacher scoring, student forward/backward,
total wall time; off-policy teacher generation runs once, outside the training loop) and peak allocated memory.}
\label{tab:t3cost}
\tab{t3_cost}
\end{table}

\begin{table}[!htb]
\centering\footnotesize
\caption{Task 3. Standalone accuracy (\%, 95\% CI, 2 samples per prompt) and truncation rate of the target and each draft
on the 60 GSM8K test prompts (up to 512 tokens) and the 60 short-form prompts (up to 32 tokens).}
\label{tab:t3acc}
\tab{t3_accuracy}
\end{table}

\section{Qualitative examples}
All examples were chosen by rules fixed in \texttt{report/make\_report\_assets.py} before their outputs were inspected,
and all use matched inputs across the compared models or policies. \emph{Task~1} (Table~\ref{tab:t1qual}): the first two
HellaSwag examples in the locked order, then the first two, in order, on which FP16 is correct and P3 ($s=0.5$) or Q3
(3 bits) chooses a different ending from FP16. \emph{Task~2C} (Tables~\ref{tab:t2cqual} and~\ref{tab:t2cqualtext}): at
4,096 tokens, the first prompt identifier for each evidence position. \emph{Task~3} (Table~\ref{tab:t3qual}): the first two
GSM8K prompts and the first short-form prompt by identifier, first repetition, for the target, the original draft and
the seed-1 off- and on-policy drafts. In Task~1 the compressed models make different mistakes on the same examples, and
Q3 even recovers example~0 where Q2 fails. In Task~2C sink+recent answers with the same code-like number (12345) on both
prompts whose evidence it evicted, and the token-wise INT4 cache produces no code. In Task~3 every draft runs past 512 tokens on
both math prompts, and all models answer the short-form prompt identically.

\begin{table}[!htb]
\centering\footnotesize
\caption{Task 1. HellaSwag cases: chosen ending (index; \checkmark\ correct) under FP16, P2/P3 at 50\% sparsity and
Q2/Q3 at 3 bits.}
\label{tab:t1qual}
\tab{t1_qual}
\end{table}

\begin{table}[!htb]
\centering\footnotesize
\caption{Task 2C. Extracted answer per cache policy for matched prompts at 4,096 tokens (\checkmark: exact match).}
\label{tab:t2cqual}
\tab{t2c_qual}
\end{table}

\begin{table}[!htb]
\centering\footnotesize
\caption{Task 2C. Raw greedy outputs for the first prompt of Table~\ref{tab:t2cqual}.}
\label{tab:t2cqualtext}
\tab{t2c_qual_text}
\end{table}

\begin{table}[!htb]
\centering\footnotesize
\caption{Task 3. Sampled standalone responses (temperature 1) for matched prompts: length, truncation at the cap,
correctness and the first 75 characters.}
\label{tab:t3qual}
\tab{t3_qual}
\end{table}

\end{document}
